% Generated draft: requires mathematical and source review.
\documentclass[11pt]{article}
\usepackage[margin=25mm]{geometry}
\usepackage{fontspec}
\setmainfont{Latin Modern Roman}
\usepackage{amsmath,amssymb,mathtools,unicode-math}
\setmathfont{Latin Modern Math}
\usepackage{enumitem,needspace,xurl,hyperref}
\setlength{\parindent}{0pt}
\setlength{\parskip}{5pt}
\setlength{\emergencystretch}{3em}
\newcommand{\N}{\mathbb{N}}
\newcommand{\Z}{\mathbb{Z}}
\newcommand{\Q}{\mathbb{Q}}
\newcommand{\R}{\mathbb{R}}
\newcommand{\C}{\mathbb{C}}
\newcommand{\F}{\mathbb{F}}
\newcommand{\A}{\mathbb{A}}
\newcommand{\PP}{\mathbb{P}}
\newcommand{\E}{\mathbb{E}}
\newcommand{\eps}{\varepsilon}
\newcommand{\dd}{\,\mathrm{d}}
\begin{document}
\section*{UnsolvedMath ID: 30003235}
\textbf{Source problem code:} catalogue entry
\subsection{Definitions and mathematical statement}
For real parameters \(a,b\), let
\[
L_{a,b}=\{(x,ax+b):x\in\R\}\subseteq\R^2.
\]
The real domain for \(a,b\) is the intended interpretation of the supplied affine-line notation; it is not explicitly stated in the record. For \(t\in\R\), define \(\|t\|=\inf_{m\in\Z}|t-m|\), the distance to the nearest integer. For positive weights \(i,j>0\) satisfying \(i+j=1\), put \(\sigma=\min\{i,j\}\) and define the standard weighted badly approximable set by
\[
\operatorname{Bad}(i,j)=\left\{(u,v)\in\R^2:\ \exists c>0\ \forall q\in\Z_{\geq1},\quad
\max\{q^i\|qu\|,q^j\|qv\|\}\geq c\right\}.
\]
Here \(\dim_H\) denotes Hausdorff dimension for the Euclidean metric. The source does not explicitly define \(\operatorname{Bad}(i,j)\); the displayed definition supplies its standard positive-weight interpretation.

The precise positive-weight question is whether, for every \(a,b\in\R\) and every \(i,j>0\) with \(i+j=1\), the implication
\[
\liminf_{\substack{q\to\infty\\q\in\Z_{\geq1}}}
q^{1/\sigma}\max\{\|qa\|,\|qb\|\}>0
\quad\Longrightarrow\quad
\dim_H\bigl(L_{a,b}\cap\operatorname{Bad}(i,j)\bigr)=1
\]
holds. The record describes the conclusion as known under the stronger hypothesis
\[
\exists\varepsilon>0:\quad
\liminf_{\substack{q\to\infty\\q\in\Z_{\geq1}}}
q^{1/\sigma-\varepsilon}\max\{\|qa\|,\|qb\|\}>0,
\]
and asks whether one may replace this hypothesis by its \(\varepsilon=0\) version. This description of the known result is supplied source information, not an independently verified theorem attribution.

The supplied weight range is actually \(i,j\geq0\). If either weight is zero, then \(\sigma=0\) and \(1/\sigma\) is undefined. The record provides neither a replacement endpoint hypothesis nor an endpoint convention for \(\operatorname{Bad}(i,j)\). Consequently the displayed implication formulates the positive-weight portion only; the requested zero-weight cases remain unresolved at the level of formulation.
\subsection{Short English statement}
For positive weights summing to one, ask whether an affine line contains a set of weighted badly approximable points of Hausdorff dimension one under the stated Diophantine condition with \(\varepsilon=0\). The source also allows zero weights, but its hypothesis is undefined there. Its open status and description of the known stronger hypothesis are supplied claims requiring review.
\subsection{Sources}
\begin{itemize}
\item Supplied local UnsolvedMath registry record, permanent ID \(30003235\), source problem code \texttt{catalogue entry}; registry locations identified by the project are \texttt{data/unsolvedmath/problems.json} and \texttt{catalog/problems.json}. Title: \emph{Weighted Bad Approximation on Affine Lines}. Collection: Oberwolfach Reports Open Problems. No individual proposer is supplied. The proposed year is recorded as 2016.
\item Supplied report citation: \emph{Arbeitsgemeinschaft: Diophantine Approximation, Fractal Geometry and Dynamics} (2017), Oberwolfach Reports, pp.\ 2749--2792, DOI 10.4171/owr/2016/48. \url{https://doi.org/10.4171/owr/2016/48}. The same URL is listed among the literature sources as ``Oberwolfach Report 48.'' The background describes these reports as records of mathematical workshops and their open-problem discussions.
\item Supplied literature source: Beresnevich--Nesharim--Yang, ``winning on curves.'' \url{https://arxiv.org/abs/2005.02128}. No fuller bibliographic citation or specific theorem is supplied.
\item Literature checked at: \texttt{2026-08-21}. This is supplied dated curation, not independently verified source checking. The supplied status is ``open.'' The assessment states that a cited Jacobian-rank theorem proves uncountability under a full-rank hypothesis, while a proposed relaxation to zero excess dimension is not established by the later curve and manifold results located. These terms do not identify the affine-line hypothesis above, and their relevance requires source review. The background explicitly calls this a dated literature triage derived from the supplied OWR curation and requests independent verification of time-sensitive status and cited claims.
\item The registry labels the statement \texttt{corrected\_verified} and reports manual semantic, typographic, and LaTeX review asserting that the clean statement is self-contained and preserves the mathematical target. Those are supplied review labels; they do not resolve the zero-weight issue or establish independent mathematical verification of this draft.
\end{itemize}
\end{document}
